\section{Portfolio and Recursive-Utility Verification}
\subsection{Stationary Merton Benchmark}
Write $b=1-\gamma$, $V(X)=AX^b/b$, $c=mX$, and
\[
 g(m,\pi)=r+(\mu-r)\pi-m+\tfrac12(b-1)\sigma^2\pi^2.
\]
For exponential utility discount $\rho$, the fixed-policy stationary equation is
\[
 A=\frac{m^b}{\rho-bg(m,\pi)}
\]
on the positive branch $\rho-bg>0$. Optimization gives $\pi^*=(\mu-r)/(\gamma\sigma^2)$ and
\[
 m^*=\frac{\rho-(1-\gamma)[r+(\mu-r)^2/(2\gamma\sigma^2)]}{\gamma}.
\]
At $(\rho,\gamma,\mu,r,\sigma)=(0.04,2,0.08,0.02,0.2)$ the targets are $(m^*,\pi^*)=(0.04125,0.75)$. Wealth remains positive under this proportional policy. Its discounted terminal candidate expectation decays at rate $\rho-bg(m^*,\pi^*)=m^*>0$, verifying its own transversality condition. The usual verification comparison is made over policies for which the stopped local martingale and discounted terminal inequalities pass to the limit.

The executed feature solver treats $q=\log A$, the consumption ratio, and the portfolio ratio as unknowns. It alternates the fixed-policy scalar evaluation root with an automatic-differentiation actor optimization, without feeding the analytical targets into the updates. Positivity is imposed by $A=e^q$. This is a learned homothetic feature model; the ratios alone would not be a neural-function benchmark on an unrestricted state space. The no-short regression at $\mu=0.01$ instead has unconstrained share $-0.125$ and constrained share zero at every positive wealth. It is not a wealth-dependent switching problem.

\subsection{Epstein--Zin Driver and the Admissible Comparison Class}
Use $a=1-1/\psi$, $b=1-\gamma$, with $a,b\ne0$, $c>0$ and $bV>0$. The driver is
\[
 f(c,V)=\frac\rho a\{c^a(bV)^{1-a/b}-bV\}.
\]
The log limits must be taken separately when a denominator vanishes. For a homothetic candidate $V=AX^b/b$, dividing the stationary HJB by $AX^b$ gives
\[
 \frac\rho a\{m^aA^{-a/b}-1\}+g(m,\pi)=0.
\]
Thus a fixed proportional policy has a unique positive homothetic coefficient whenever $1-ag/\rho>0$. Its consumption first-order condition is $\rho m^{a-1}A^{-a/b}=1$, and the portfolio first-order condition gives the same Merton share with risk aversion $\gamma$. Eliminating $A$ yields
\begin{align*}
 \pi^*&=\frac{\mu-r}{\gamma\sigma^2},\\
 m^*&=\psi\rho+(1-\psi)\left[r+\frac{(\mu-r)^2}{2\gamma\sigma^2}\right],\\
 A^*&=(\rho(m^*)^{a-1})^{b/a}.
\end{align*}
For $(\rho,\gamma,\psi,\mu,r,\sigma)=(0.04,5,1.5,0.08,0.02,0.2)$ this gives $m^*=0.0455$ and $\pi^*=0.3$. The normalized Hamiltonian is strictly concave in the consumption ratio and portfolio share on the specified positive domain. The selected ratios lie inside the implemented action bounds.

For precision about infinite horizon, the verification class consists of policies for which the policy-wise recursive equation exists on the branch $bY>0$, its localized comparison argument is integrable, and the comparison-weighted terminal remainder vanishes. On compact subsets of $c>0,bY>0$ the driver is locally Lipschitz, but this does not imply a global Lipschitz constant on the entire half-line. One may first stop at compact localization sets, apply the mean-value linearization in the value argument, and then impose uniform integrability and the weighted terminal limit to remove localization. With $\Gamma_t=\exp(\int_0^t f_V(c_s,\zeta_s)ds)$ for the intermediate value $\zeta_s$, the required remainder is of the form $\E[\Gamma_T|V(X_T)-Y_T^p|]\to0$. The candidate HJB inequality and this comparison yield $Y_0^p\leq V(X_0)$, with equality for the maximizing policy.

For the selected homothetic policy itself, $f_V=-0.1115$, while the growth exponent of $\E X_t^b$ is $bg=0.066$. The candidate comparison-weighted value therefore decays at rate $0.1115-0.066=0.0455>0$. The numerical record evaluates the fixed-policy equation and target errors, and the analytic calculation verifies this decay. This specifies a genuine stationary recursive benchmark with an explicit admissible verification class; it does not assert that all bounded controls satisfy the terminal condition automatically.

The sign-enforcing architecture is actually executed: $V=e^qX^b/b$, hence $bV=e^qX^b>0$. A separate finite-horizon architecture test uses
\[
 V_\xi(t,X)=\frac1b\exp\{\log[bG(X)]+(T-t)N_\xi(t,X)\}
\]
for $bG>0$. It verifies the terminal identity and sign on test points. It is a domain test, not evidence that an additional finite-horizon recursive problem was trained. These distinctions follow the policy-wise recursive-utility formulation of \citet{duffie1992} rather than replacing recursion by a pointwise evaluation at $c=1,V=-1$.

\subsection{A Structured Stochastic Dynamic Control Check}
For $d\in\{4,8,16\}$, the additional six-period LQR problem uses
\[
 A=0.82I+0.07(J+J^\top),\quad B=0.2I,\quad
 Q=I+0.3\mathbf1\mathbf1^\top/d,\quad R=0.15I,
\]
and Gaussian shock loading $\Sigma=0.08(I+0.3\mathbf1\mathbf1^\top/d)$, with terminal matrix $Q$. Here $J$ has ones on its upper adjacent diagonal. The linear neural actor and quadratic feature critic are trained by policy evaluation and improvement. The independent reference uses the finite-horizon Riccati recursion, including the noise contribution to the value constant. A second recursion evaluates the learned fixed policy rather than using its fitted critic as a performance score.

\begin{table}[tbp]
\centering
\caption{Structured Stochastic LQR Regression}
\label{tab:lqr}
\small
\begin{tabular}{@{}lrrrr@{}}
\toprule
$d$ & Actor coeff. error & Cost excess max. & Neural sec. & Riccati sec. \\
\midrule
4 & 4.25e-06 & 3.03e-11 & 0.0571 & 0.00037 \\
8 & 2.25e-06 & 2.14e-11 & 0.0724 & 0.00028 \\
16 & 4.30e-06 & 3.84e-11 & 0.1330 & 0.00033 \\
\bottomrule
\end{tabular}
\par\smallskip\begin{minipage}{\linewidth}\footnotesize Six periods, linear actors and quadratic feature critics, independent Riccati reference, and separate frozen-policy evaluation. Cost excess is normalized per dimension in the implementation. This structured unconstrained regression is not the main binding-resource experiment and supplies no generic neural speed advantage.\end{minipage}
\end{table}


The experiment has cross-state dynamics and non-diagonal stochastic covariance but no binding control constraints. It is retained as a structured dynamic regression test. The main shared-budget resource experiment supplies the separately executed constrained case; neither is relabeled as an adaptive sparse-grid comparison.
